\chapter{Realized Volatility}
\label{ch:rv}

\begin{application}[Why This Chapter]
Realized volatility is the target variable for every project direction in this guide.
Chapters~\ref{ch:noise} and \ref{ch:jumps} refine the estimator (handling noise and jumps), Chapter~\ref{ch:har} builds the HAR forecasting model around it, and Chapter~\ref{ch:features} uses RV-derived features.
You need to understand RV at an intuitive level (what it measures, how it is constructed, why 5-minute sampling) before any of that.
\end{application}

Chapter~\ref{ch:returns} introduced variance as a measure of how spread out returns are.
That chapter computed variance from daily returns over weeks or months, producing a single number for the whole sample: unconditional volatility.
This chapter moves to a different question: how volatile was the market \emph{today}?
Answering that requires looking inside the trading day at high-frequency (intraday) returns.

% ═══════════════════════════════════════════════════
\section{The Theoretical Target: Integrated Variance}
\label{sec:integrated-variance}
% ═══════════════════════════════════════════════════

Before building an estimator, you need to know what you are estimating.
The theoretical quantity is called \emph{integrated variance}.

\begin{prereq}[Integrals as Accumulated Area]
If $f(x)$ is a non-negative function, the integral $\int_a^b f(x)\,dx$ equals the area under the curve $f$ between $a$ and $b$.
When $f$ varies over time, the integral accumulates all of those local values into a single total.
You can think of $\int_a^b f(x)\,dx$ as ``add up all the tiny contributions of $f$ across the interval $[a,b]$.''
\end{prereq}

\begin{prereq}[Price Processes (Informal)]
In continuous-time finance, a price process $P_t$ is modeled as a \emph{semimartingale}.
For this chapter, the key property is that the log price $p_t = \ln P_t$ evolves according to
\[
  dp_t = \mu_t\,dt + \sigma_t\,dW_t
\]
where $\mu_t$ is the drift (expected return per unit time), $\sigma_t$ is the instantaneous volatility (how rapidly the price fluctuates at time $t$), and $W_t$ is a Wiener process (continuous random walk).
The drift is usually negligible at intraday horizons; the volatility $\sigma_t$ is the quantity we care about.
You do not need to understand stochastic calculus to follow this chapter.
The key takeaway: $\sigma_t$ can change from moment to moment within the day.
\end{prereq}

The intuition is straightforward.
At every instant during the trading day, there is a local volatility level $\sigma_t$ describing how rapidly the price fluctuates.
This level is not constant; it changes as news arrives, liquidity shifts, and traders adjust positions.
\emph{Integrated variance} adds up all of these instantaneous variance levels across the day to produce a single summary of total price variation.

Think of it by analogy: if you drive a car at varying speeds throughout the day, the total distance driven is the integral of your speed over time.
Integrated variance is the ``total distance'' of price fluctuation.

\begin{definition}[Integrated Variance]
For a log-price process $p_t$ with instantaneous variance $\sigma^2_t$, the integrated variance over day $t$ (from the close of day $t-1$ to the close of day $t$) is:
\begin{equation}
  \IVol_t = \int_{t-1}^{t} \sigma^2_s \, ds
  \label{eq:iv}
\end{equation}
\begin{itemize}[nosep]
  \item $\IVol_t$: integrated variance over day $t$
  \item $\sigma^2_s$: instantaneous variance at time $s$ within the day
  \item $ds$: an infinitesimal increment of time
  \item The integral sums the instantaneous variances from the previous close ($t-1$) to today's close ($t$)
\end{itemize}
\end{definition}

Integrated variance is a \emph{latent} quantity: you never observe $\sigma^2_s$ directly.
The entire point of this chapter is to estimate $\IVol_t$ from observable intraday price data.

\begin{keyidea}[Integrated Variance = Energy of Price Movements]
$\IVol_t$ measures the total ``energy'' of price movements during day $t$.
A day with high $\IVol_t$ is a day where the price path wiggled a lot (even if the open-to-close return was small).
A day with low $\IVol_t$ is a day where the price drifted smoothly.
\end{keyidea}

% ═══════════════════════════════════════════════════
\section{Realized Variance as an Estimator}
\label{sec:rv-estimator}
% ═══════════════════════════════════════════════════

You now know the target: integrated variance $\IVol_t$.
The question is how to estimate it from data you can actually observe.
The answer, developed by \citet{ABDL2001} and \citet{BNS2002}, is remarkably simple: sum up squared intraday returns.

\subsection{Construction}

Divide the trading day into $n$ equal intervals.
At the end of each interval, record the price.
Compute the log return over each interval.
Square each return and sum.

\begin{definition}[Realized Variance]
Given $n$ intraday log returns $r_{t,1}, r_{t,2}, \ldots, r_{t,n}$ within day $t$, the realized variance is:
\begin{equation}
  \RV_t = \sum_{i=1}^{n} r^2_{t,i}
  \label{eq:rv}
\end{equation}
\begin{itemize}[nosep]
  \item $\RV_t$: realized variance for day $t$
  \item $r_{t,i}$: the $i$-th intraday log return on day $t$; that is, $r_{t,i} = p_{t,i} - p_{t,i-1}$, where $p_{t,i}$ is the log price at the end of the $i$-th interval
  \item $n$: the number of intraday intervals (e.g., $n = 78$ for 5-minute intervals over a 6.5-hour U.S.\ equity trading day)
  \item No mean subtraction: the sample mean of intraday returns is so close to zero that omitting it improves finite-sample performance \citep{ABDL2003}
\end{itemize}
\end{definition}

\begin{intuition}[Why Squaring Works]
Consider a single intraday log return $r_{t,i} = \sigma_{t,i}\,\epsilon_i$, where $\sigma_{t,i}$ is the local volatility level and $\epsilon_i$ is a noise term with mean zero and variance one.
Squaring gives $r^2_{t,i} \approx \sigma^2_{t,i}\,\epsilon^2_i$.
Since $\E[\epsilon^2_i] = 1$, each squared return is (on average) an estimate of the local variance.
Summing them across the day accumulates these local estimates into a total.
As the intervals shrink ($n \to \infty$), the noise in each $\epsilon^2_i$ averages out, and the sum converges to the integrated variance.
\end{intuition}

\subsection{Why It Works: Convergence to Quadratic Variation}

The theoretical justification comes from the theory of \emph{quadratic variation}.

\begin{prereq}[Quadratic Variation (Informal)]
For a continuous-time stochastic process $X_t$, the quadratic variation $[X]_t$ measures the cumulative squared increments of the path up to time $t$.
It is defined as the limit of the sum of squared increments as the partition becomes infinitely fine:
\[
  [X]_t = \lim_{n \to \infty} \sum_{i=1}^{n} (X_{t_i} - X_{t_{i-1}})^2
\]
For a process with continuous paths (no jumps), the quadratic variation equals the integrated variance: $[X]_t = \int_0^t \sigma^2_s\,ds$.
If the process has jumps, the quadratic variation also captures the sum of squared jumps.
\end{prereq}

As you sample more frequently (as $n \to \infty$ and the length of each interval $\Delta \to 0$), $\RV_t$ converges to the quadratic variation of the log-price process \citep{ABDL2001, BNS2002}:

\begin{equation}
  \RV_t \xrightarrow{p} [p]_t \quad \text{as } \Delta \to 0
  \label{eq:rv-convergence}
\end{equation}
\begin{itemize}[nosep]
  \item $\xrightarrow{p}$: convergence in probability
  \item $[p]_t$: quadratic variation of the log-price process on day $t$
  \item $\Delta = 1/n$: the length of each sampling interval (as a fraction of the trading day)
\end{itemize}

If the price path is continuous (no jumps), then the quadratic variation equals the integrated variance:
\begin{equation}
  [p]_t = \IVol_t \qquad (\text{no jumps})
  \label{eq:qv-equals-iv}
\end{equation}

If the price path has jumps, the quadratic variation includes a jump component:
\begin{equation}
  [p]_t = \IVol_t + \sum_{s \leq t} (\Delta p_s)^2 \qquad (\text{with jumps})
  \label{eq:qv-with-jumps}
\end{equation}
\begin{itemize}[nosep]
  \item $\Delta p_s$: the jump in log price at time $s$ (zero if no jump occurs at $s$)
  \item The sum collects all jumps within the day
  \item Separating jumps from continuous variation is the subject of Chapter~\ref{ch:jumps}
\end{itemize}

\begin{keyresult}[\citet{ABDL2001}, \citet{BNS2002}: RV Consistency]
Under mild regularity conditions, realized variance $\RV_t$ is a consistent estimator of quadratic variation $[p]_t$.
In the absence of jumps, $\RV_t$ consistently estimates integrated variance $\IVol_t$.
This result holds regardless of the specific form of $\sigma_t$; the volatility process can be stochastic, path-dependent, or driven by latent factors.
\end{keyresult}

\subsection{Diagram: Building RV from a Price Path}

Figure~\ref{fig:rv-construction} shows how RV is constructed from a single day's price data.

\begin{figure}[htbp]
\centering
\begin{tikzpicture}
\begin{axis}[
  width=\textwidth,
  height=7cm,
  xlabel={Time of day},
  ylabel={Price (\$)},
  xmin=0, xmax=6.5,
  ymin=99.0, ymax=101.5,
  xtick={0, 0.5, 1, 1.5, 2, 2.5, 3, 3.5, 4, 4.5, 5, 5.5, 6, 6.5},
  xticklabels={9:30, 10:00, 10:30, 11:00, 11:30, 12:00, 12:30, 1:00, 1:30, 2:00, 2:30, 3:00, 3:30, 4:00},
  xticklabel style={rotate=45, anchor=east, font=\scriptsize},
  grid=major,
  grid style={gray!20},
  tick label style={font=\small},
  label style={font=\small},
  title style={font=\small\bfseries},
  title={Constructing RV: Price Path Sampled at 30-Minute Intervals},
]
% Continuous price path (smooth interpolation)
\addplot[gray!50, thick, mark=none, smooth] coordinates {
  (0, 100.00) (0.25, 100.15) (0.5, 100.30) (0.75, 100.22)
  (1.0, 100.45) (1.25, 100.55) (1.5, 100.60) (1.75, 100.50)
  (2.0, 100.35) (2.25, 100.20) (2.5, 100.10) (2.75, 100.25)
  (3.0, 100.40) (3.25, 100.55) (3.5, 100.70) (3.75, 100.60)
  (4.0, 100.50) (4.25, 100.42) (4.5, 100.55) (4.75, 100.68)
  (5.0, 100.80) (5.25, 100.75) (5.5, 100.65) (5.75, 100.58)
  (6.0, 100.70) (6.25, 100.78) (6.5, 100.85)
};

% Sampled points (half-hourly)
\addplot[only marks, mark=*, mark size=3pt, defblue] coordinates {
  (0, 100.00) (0.5, 100.30) (1.0, 100.45)
  (1.5, 100.60) (2.0, 100.35) (2.5, 100.10)
  (3.0, 100.40) (3.5, 100.70) (4.0, 100.50)
  (4.5, 100.55) (5.0, 100.80) (5.5, 100.65)
  (6.0, 100.70) (6.5, 100.85)
};

% Connect sampled points
\addplot[defblue, thick, mark=none] coordinates {
  (0, 100.00) (0.5, 100.30) (1.0, 100.45)
  (1.5, 100.60) (2.0, 100.35) (2.5, 100.10)
  (3.0, 100.40) (3.5, 100.70) (4.0, 100.50)
  (4.5, 100.55) (5.0, 100.80) (5.5, 100.65)
  (6.0, 100.70) (6.5, 100.85)
};

% Annotations for a few returns
\draw[<->, warnred, thick] (axis cs: 0, 99.3) -- (axis cs: 0.5, 99.3)
  node[midway, below, font=\scriptsize] {$r_{t,1}$};
\draw[<->, warnred, thick] (axis cs: 0.5, 99.3) -- (axis cs: 1.0, 99.3)
  node[midway, below, font=\scriptsize] {$r_{t,2}$};
\draw[<->, warnred, thick] (axis cs: 1.0, 99.3) -- (axis cs: 1.5, 99.3)
  node[midway, below, font=\scriptsize] {$r_{t,3}$};

% RV formula callout
\node[anchor=west, font=\small, fill=white, draw=defblue, rounded corners, inner sep=4pt]
  at (axis cs: 3.5, 99.5) {$\RV_t = r_{t,1}^2 + r_{t,2}^2 + \cdots + r_{t,13}^2$};

\end{axis}
\end{tikzpicture}
\caption{Constructing realized variance from a price path.
Gray: the continuous intraday price.
Blue dots: prices sampled at half-hourly intervals ($n = 13$ intervals).
Each interval produces a log return $r_{t,i}$.
RV is the sum of all squared returns.
Sampling more frequently (e.g., every 5 minutes instead of 30) produces more terms in the sum and a more precise estimate.}
\label{fig:rv-construction}
\end{figure}

\begin{workedexample}{Computing RV from Half-Hourly Returns}
A stock opens at \$100.00 and you record its price every 30 minutes over a shortened trading day (6 intervals for simplicity).
The prices and resulting log returns are:

\medskip
\begin{center}
\begin{tabular}{@{} c c c c @{}}
\toprule
\textbf{Time} & \textbf{Price (\$)} & \textbf{Log return $r_{t,i}$} & \textbf{$r^2_{t,i}$} \\
\midrule
9:30  (open)  & 100.00 &       &       \\
10:00         & 100.30 & $\ln(100.30/100.00) = +0.003000$ & $9.000 \times 10^{-6}$ \\
10:30         & 100.10 & $\ln(100.10/100.30) = -0.001998$ & $3.992 \times 10^{-6}$ \\
11:00         & 100.50 & $\ln(100.50/100.10) = +0.003992$ & $15.936 \times 10^{-6}$ \\
11:30         & 100.20 & $\ln(100.20/100.50) = -0.002992$ & $8.952 \times 10^{-6}$ \\
12:00         & 100.60 & $\ln(100.60/100.20) = +0.003988$ & $15.904 \times 10^{-6}$ \\
12:30 (close) & 100.45 & $\ln(100.45/100.60) = -0.001491$ & $2.223 \times 10^{-6}$ \\
\bottomrule
\end{tabular}
\end{center}
\medskip

\textbf{Step 1: Sum the squared returns.}
\[
  \RV_t = (9.000 + 3.992 + 15.936 + 8.952 + 15.904 + 2.223) \times 10^{-6} = 56.007 \times 10^{-6}
\]

\textbf{Step 2: Convert to realized volatility.}
\[
  \sqrt{\RV_t} = \sqrt{56.007 \times 10^{-6}} = 0.00748 \quad (0.75\% \text{ daily})
\]

\textbf{Step 3: Annualize.}
\[
  \text{Annualized realized volatility} = 0.00748 \times \sqrt{252} = 0.119 \quad (11.9\%)
\]

For context, 11.9\% annualized is a bit below the S\&P~500's long-run average (roughly 15--20\%).
With only 6 intervals, this estimate is noisy.
Sampling every 5 minutes would give 78 intervals over a 6.5-hour day, producing a much more precise estimate.
\end{workedexample}

% ═══════════════════════════════════════════════════
\section{How Frequently to Sample}
\label{sec:sampling-frequency}
% ═══════════════════════════════════════════════════

You now have an estimator ($\RV_t$) and a convergence result (it approaches $[p]_t$ as $n \to \infty$).
The obvious next step is to sample as frequently as possible: every second, every tick, every trade.
In theory, this is optimal.
In practice, it fails.

\subsection{The Theory: More Is Better}

The convergence result in Equation~\eqref{eq:rv-convergence} says that RV becomes a better estimator as $\Delta$ shrinks.
More precisely, \citet{BNS2002} showed that the estimation error has a central limit theorem:

\begin{equation}
  \sqrt{n}\,(\RV_t - \IVol_t) \xrightarrow{d} \N\!\left(0,\; 2\int_{t-1}^{t}\sigma^4_s\,ds\right)
  \label{eq:rv-clt}
\end{equation}
\begin{itemize}[nosep]
  \item $\xrightarrow{d}$: convergence in distribution
  \item $\sqrt{n}$: the estimation error shrinks at rate $1/\sqrt{n}$
  \item $2\int_{t-1}^{t}\sigma^4_s\,ds$: the asymptotic variance (it depends on the ``quarticity,'' the integral of $\sigma^4$)
  \item The message: doubling the number of intervals cuts the standard error by a factor of $\sqrt{2}$
\end{itemize}

So with perfect data, you would sample every millisecond and get an essentially noise-free estimate of integrated variance.

\subsection{The Practice: Microstructure Noise}

Real transaction prices are not clean readings of a frictionless price process.
They are contaminated by \emph{microstructure noise}: the cumulative effect of bid-ask bounce, discrete price grids, order-processing delays, and other trading frictions.

\begin{prereq}[Bid-Ask Bounce]
When you buy a stock, you pay the \emph{ask} price (slightly above the ``true'' value).
When you sell, you receive the \emph{bid} price (slightly below).
The difference between bid and ask is the \emph{spread}.
Even if the true value is perfectly constant, alternating buys and sells cause the transaction price to bounce between bid and ask.
This artificial oscillation creates spurious volatility in very high-frequency returns.
\end{prereq}

The standard model for microstructure noise \citep{ABDL2001, BNS2002} adds a noise term to the observed log price:
\begin{equation}
  p^*_{t,i} = p_{t,i} + \varepsilon_{t,i}
  \label{eq:noise-model}
\end{equation}
\begin{itemize}[nosep]
  \item $p^*_{t,i}$: observed (noisy) log price
  \item $p_{t,i}$: true (latent) efficient log price
  \item $\varepsilon_{t,i}$: microstructure noise, typically assumed i.i.d.\ with $\E[\varepsilon_{t,i}] = 0$ and $\operatorname{Var}(\varepsilon_{t,i}) = \omega^2$
\end{itemize}

When you compute returns from noisy prices and sample very frequently, the noise dominates.
Instead of converging to $\IVol_t$, realized variance computed from noisy prices diverges:
\begin{equation}
  \RV_t^{(\text{noisy})} \to 2n\omega^2 \quad \text{as } n \to \infty
  \label{eq:rv-noise-divergence}
\end{equation}

The estimate explodes linearly in $n$ because each noisy return contributes roughly $2\omega^2$ on average, and you are summing $n$ of them.

\begin{warning}[Sampling Too Fast Destroys the Estimate]
With tick-by-tick or second-by-second data, the bid-ask bounce inflates realized variance far above the true integrated variance.
Sampling more frequently makes the estimate \emph{worse}, not better.
This is the fundamental tension of realized variance estimation.
\end{warning}

\subsection{The Bias-Variance Tradeoff}

The result is a tradeoff:
\begin{itemize}[nosep]
  \item \textbf{Sample too infrequently} (e.g., hourly): few observations, high estimation variance, but little noise contamination.
  \item \textbf{Sample too frequently} (e.g., every second): many observations, low estimation variance in theory, but massive upward bias from microstructure noise.
\end{itemize}

There is an optimal sampling frequency that balances these two forces.
The concept is often visualized with a \emph{volatility signature plot}.

\subsection{The Volatility Signature Plot}

\begin{definition}[Volatility Signature Plot]
A volatility signature plot displays the average realized variance (or realized volatility), computed across many days, as a function of the sampling frequency.
The $x$-axis shows the return interval (e.g., 1 second, 1 minute, 5 minutes, \ldots), and the $y$-axis shows the resulting average $\RV_t$.
\end{definition}

Figure~\ref{fig:vol-sig-plot} shows the typical shape.

\begin{figure}[htbp]
\centering
\begin{tikzpicture}
\begin{axis}[
  width=0.85\textwidth,
  height=7cm,
  xlabel={Sampling interval},
  ylabel={Average realized variance ($\times 10^{-4}$)},
  xmin=0, xmax=11,
  ymin=0.5, ymax=4.5,
  xtick={1, 2, 3, 4, 5, 6, 7, 8, 9, 10},
  xticklabels={1s, 5s, 15s, 30s, 1m, 5m, 15m, 30m, 1h, 2h},
  grid=major,
  grid style={gray!20},
  tick label style={font=\small},
  label style={font=\small},
  title style={font=\small\bfseries},
  title={Volatility Signature Plot (Schematic)},
]

% Noisy region: RV inflated at high frequencies
\addplot[warnred, very thick, mark=none, smooth] coordinates {
  (1, 4.2) (2, 3.1) (3, 2.2) (4, 1.7) (5, 1.35)
};

% Stable region: RV approximately flat
\addplot[intgreen, very thick, mark=none, smooth] coordinates {
  (5, 1.35) (6, 1.15) (7, 1.10) (8, 1.08)
};

% Low-frequency region: increased estimation variance
\addplot[keyorange, very thick, mark=none, smooth] coordinates {
  (8, 1.08) (9, 0.95) (10, 0.75)
};

% True IV level
\addplot[dashed, defblue, thick, mark=none] coordinates {(0, 1.1) (11, 1.1)};

% Annotations
\node[anchor=south, font=\scriptsize, warnred] at (axis cs: 2.5, 3.3) {Noise inflates RV};
\draw[-{Stealth}, warnred] (axis cs: 2.5, 3.2) -- (axis cs: 2.0, 2.8);

\node[anchor=west, font=\scriptsize, intgreen] at (axis cs: 6.2, 1.5) {``Sweet spot''};
\draw[-{Stealth}, intgreen] (axis cs: 6.2, 1.45) -- (axis cs: 6.0, 1.2);

\node[anchor=north, font=\scriptsize, keyorange] at (axis cs: 9.5, 0.65) {Few obs, high};
\node[anchor=north, font=\scriptsize, keyorange] at (axis cs: 9.5, 0.50) {estimation error};

\node[anchor=west, font=\scriptsize, defblue] at (axis cs: 7.5, 1.22) {True $\IVol_t$};

\end{axis}
\end{tikzpicture}
\caption{Schematic volatility signature plot.
At very high frequencies (left, red), microstructure noise inflates realized variance.
At very low frequencies (right, orange), too few observations make the estimate imprecise and biased downward.
The green region (roughly 1--15 minutes) is the ``sweet spot'' where RV is closest to the true integrated variance (dashed blue line).
Chapter~\ref{ch:noise} develops noise-robust estimators that work at higher frequencies.}
\label{fig:vol-sig-plot}
\end{figure}

The plot has three regions:
\begin{enumerate}[nosep]
  \item \textbf{High-frequency region} (left): RV rises sharply, contaminated by bid-ask bounce and microstructure effects.
  \item \textbf{Intermediate region} (center): RV flattens near the true integrated variance. This is the practical sweet spot.
  \item \textbf{Low-frequency region} (right): RV drifts downward and becomes noisy because too few observations produce imprecise estimates.
\end{enumerate}

\begin{keyidea}[The Fundamental Tradeoff]
In theory, sample as fast as possible.
In practice, noise wins beyond about 1-minute returns for most liquid assets.
The volatility signature plot is a diagnostic tool: if average RV is still rising as you increase frequency, you are in the noise-contaminated region and need to back off (or use a noise-robust estimator from Chapter~\ref{ch:noise}).
\end{keyidea}

% ═══════════════════════════════════════════════════
\section{5-Minute RV: The Practical Workhorse}
\label{sec:five-minute-rv}
% ═══════════════════════════════════════════════════

Given the tradeoff from the previous section, what sampling frequency should you use in practice?
The answer, for most applications, is 5 minutes.

This is not a number derived from first principles (the optimal frequency depends on the noise-to-signal ratio, which varies by asset and time period).
It is an empirical finding, and a remarkably robust one.

\begin{keyresult}[\citet{LPS2015}: Does Anything Beat 5-Minute RV?]
\citet{LPS2015} conducted a large-scale comparison of approximately 400 realized volatility estimators (including subsampled, kernel-based, pre-averaged, and multi-scale estimators) applied to 31 assets across 5 asset classes (equities, equity indices, exchange rates, bonds, and commodities).
Their conclusion: for the purpose of \emph{forecasting} future volatility, ``it is difficult to significantly outperform'' the simple 5-minute realized variance.
More sophisticated estimators sometimes produce marginally more accurate daily estimates, but these gains do not reliably translate into better forecasts.
\end{keyresult}

Why 5 minutes?
Two practical reasons:
\begin{enumerate}[nosep]
  \item For liquid assets (major equity indices, G10 FX pairs, actively traded stocks), the volatility signature plot typically flattens by the 5-minute mark. At this frequency, microstructure noise is small relative to genuine price variation.
  \item 5-minute sampling produces 78 intraday observations per day for U.S.\ equities (6.5 hours $\times$ 12 intervals per hour), which provides enough observations for the central limit theorem in Equation~\eqref{eq:rv-clt} to give a reasonable approximation.
\end{enumerate}

\begin{keyidea}[The Default Choice]
Unless you have a specific reason to do otherwise (illiquid assets, tick-level analysis, or a noise-robust estimator from Chapter~\ref{ch:noise}), use 5-minute returns to compute RV.
This is the baseline in most empirical volatility research and the starting point for every forecasting model in this guide.
\end{keyidea}

\begin{warning}[5 Minutes Is Not Universal]
For illiquid assets (small-cap stocks, emerging-market currencies, less-traded commodities), the noise-contaminated region can extend to 15 or even 30 minutes.
Always check the volatility signature plot for your specific asset before committing to a sampling frequency.
\end{warning}

% ═══════════════════════════════════════════════════
\section{Realized Volatility vs.\ Realized Variance}
\label{sec:rv-vs-rvol}
% ═══════════════════════════════════════════════════

This chapter has used ``realized variance'' for $\RV_t = \sum r^2_{t,i}$ and ``realized volatility'' for $\sqrt{\RV_t}$.
Not every paper follows this convention.

\begin{warning}[Variance vs.\ Volatility: A Factor-of-10 Error]
The S\&P~500's typical daily realized variance is around $1 \times 10^{-4}$.
Its typical daily realized volatility is $\sqrt{1 \times 10^{-4}} = 0.01$ (1\%).
Confusing the two is a factor-of-100 error in the daily number, which becomes roughly a factor of 10 when annualized ($0.01 \times \sqrt{252} \approx 0.16$ vs.\ $0.0001 \times \sqrt{252} \approx 0.0016$).
Whenever you read a paper, check whether the authors report $\RV_t$ or $\sqrt{\RV_t}$, and note the units (decimal, percent, or basis points).
\end{warning}

\begin{definition}[Conventions in This Guide]
Throughout this guide:
\begin{itemize}[nosep]
  \item $\RV_t$ (and ``RV'') always refers to \emph{realized variance}: $\RV_t = \sum_{i=1}^{n} r^2_{t,i}$
  \item Realized \emph{volatility} refers to $\sqrt{\RV_t}$, the square root of realized variance
  \item When annualizing, multiply $\sqrt{\RV_t}$ by $\sqrt{252}$ (as in Section~\ref{sec:annualizing} of Chapter~\ref{ch:returns})
\end{itemize}
These conventions follow \citet{ABDL2003} and \citet{BNS2002}.
\end{definition}

Some papers (particularly in the HAR literature, Chapter~\ref{ch:har}) work with the natural logarithm of realized variance, $\ln(\RV_t)$.
The log transform has two practical benefits:
\begin{enumerate}[nosep]
  \item $\RV_t$ is right-skewed and bounded below by zero; $\ln(\RV_t)$ is approximately Gaussian \citep{ABDL2001, ABDL2003}, making standard regression assumptions more plausible.
  \item Forecasting $\ln(\RV_t)$ automatically ensures that predictions of the variance level are positive (since $e^x > 0$ for all $x$).
\end{enumerate}
When you encounter $\ln(\RV_t)$ in later chapters, this is why.

% ═══════════════════════════════════════════════════
% SUMMARY
% ═══════════════════════════════════════════════════

\section*{Summary}
\addcontentsline{toc}{section}{Summary}

\begin{itemize}[nosep]
  \item \textbf{Integrated variance} $\IVol_t = \int_{t-1}^{t} \sigma^2_s\,ds$ is the theoretical target: the total variance accumulated within a single day.

  \item \textbf{Realized variance} $\RV_t = \sum_{i=1}^{n} r^2_{t,i}$ is a model-free estimator of integrated variance, constructed by summing squared intraday returns.

  \item \textbf{Convergence}: as the sampling frequency increases ($n \to \infty$), RV converges in probability to the quadratic variation of the log-price process \citep{ABDL2001, BNS2002}.

  \item In the \textbf{absence of jumps}, quadratic variation equals integrated variance, so RV consistently estimates $\IVol_t$.

  \item With \textbf{jumps}, quadratic variation equals integrated variance plus the sum of squared jumps. Chapter~\ref{ch:jumps} addresses separating the two components.

  \item \textbf{Microstructure noise} (bid-ask bounce, discrete prices) corrupts very high-frequency returns, causing RV to diverge as sampling frequency increases.

  \item The \textbf{volatility signature plot} (average RV vs.\ sampling interval) visualizes the bias-variance tradeoff: too fast means noise contamination, too slow means imprecise estimates.

  \item \textbf{5-minute RV} is the standard benchmark. \citet{LPS2015} showed it is hard to beat for forecasting across 31 assets and ${\sim}400$ competing estimators.

  \item For \textbf{illiquid assets}, 5 minutes may still be in the noise-contaminated region; always check the volatility signature plot.

  \item \textbf{Convention}: this guide uses $\RV_t$ for realized variance and $\sqrt{\RV_t}$ for realized volatility. Confusing the two is roughly a factor-of-10 annualized error.

  \item The \textbf{log transform} $\ln(\RV_t)$ is approximately Gaussian and guarantees positive forecasts, which is why it appears frequently in later chapters.

  \item Chapter~\ref{ch:noise} develops noise-robust estimators that can safely use higher frequencies; Chapter~\ref{ch:jumps} separates continuous and jump variation.
\end{itemize}

% ═══════════════════════════════════════════════════
% KEY RESULTS TABLE
% ═══════════════════════════════════════════════════

\begin{table}[htbp]
\centering
\caption{Key results referenced in this chapter.}
\label{tab:ch02-key-results}
\begin{tabular}{@{} p{4.5cm} p{6cm} p{4cm} @{}}
\toprule
\textbf{Paper} & \textbf{Result} & \textbf{Relevance} \\
\midrule
\citet{ABDL2001} &
Realized variance converges to quadratic variation as sampling frequency increases; the distribution of $\ln(\RV_t)$ is approximately Gaussian. &
Foundational result establishing RV as a consistent, model-free volatility estimator. \\[6pt]

\citet{ABDL2003} &
Formal econometric theory for realized variance: consistency, asymptotic distribution, practical implementation with 5-minute FX returns, and the result that mean subtraction is unnecessary. &
Provided the CLT for RV and demonstrated that dropping the intraday mean improves finite-sample accuracy. \\[6pt]

\citet{BNS2002} &
Parallel asymptotic theory for realized variance; introduced bipower variation (BPV) to separate jumps from continuous variation. &
Second foundational contribution; BPV is the basis for jump detection in Chapter~\ref{ch:jumps}. \\[6pt]

\citet{LPS2015} &
Compared ${\sim}400$ realized estimators across 31 assets in 5 asset classes; simple 5-minute RV is very hard to beat for volatility forecasting. &
Justifies 5-minute RV as the default benchmark for all forecasting models in this guide. \\
\bottomrule
\end{tabular}
\end{table}
